% ============================================================
%  DysfluentNet — Interspeech 2026 Camera-Ready Submission
%  Paper #696
%  Deadline: 2026-06-19 (AoE)
%
%  Changes from revised review version → camera-ready:
%    CR-1  \documentclass[cameraready]{Interspeech}  (already set)
%    CR-2  All blue/coloured template text removed; all text is black
%    CR-3  \usepackage{color} and \usepackage{comment} removed
%           (template blue/red macros not used in this paper)
%    CR-4  \usepackage{xspace} removed (not needed; no \xspace calls remain)
%    CR-5  \usepackage{siunitx} removed (no \SI commands used)
%    CR-6  Generative AI Use Disclosure section made mandatory and
%           placed correctly (pages 5–6 zone, after Acknowledgments)
%    CR-7  \ifcameraready guard removed from Acknowledgments so it
%           renders unconditionally in the camera-ready build
%    CR-8  All ISSUE/ISSN comment lines stripped (internal notes)
%    CR-9  Double-dash em-dashes (---) standardised throughout
%    CR-10 Abstract: LaTeX math ($...$) preserved (allowed in PDF;
%           submission form copy must be plain text — note for authors)
%    CR-11 \url{} links now reveal identity (permitted in camera-ready)
%    CR-12 "In our previous work" phrasing added where self-reference
%           was anonymised (none required here — third-person already used)
%    CR-13 No technical content beyond page 4 (acknowledgments,
%           AI disclosure, and references fill pages 5–6)
% ============================================================
\documentclass[cameraready]{Interspeech}

%------------------------------------------------------------------
% Title — must match CMT submission exactly
%------------------------------------------------------------------
\title{DysfluentNet: Joint Stuttering Event Detection and Dysfluency-Aware
       Transcription via Hierarchical Self-Supervised Learning}

%------------------------------------------------------------------
% Authors — order must match CMT submission
%------------------------------------------------------------------
\author[affiliation={1},orcid=0009-0000-1917-2732, equalcontribution]{Mohankumar}{Muthu}
\author[affiliation={1},orcid=0000-0002-1848-7692, equalcontribution, correspondingauthor]{Sasikala}{E}
\author[affiliation={2},orcid=0000-0003-4978-542X]{Girirajan}{S}

\address{
$^1$ Department of Data Science and Business Systems, SRM Institute of Science and Technology, Kattankulathur--603203, Chennai, India \\
$^2$ Department of Computing Technologies, SRM Institute of Science and Technology, Kattankulathur--603203, Chennai, India
}

%------------------------------------------------------------------
% Contact
%------------------------------------------------------------------
\email{mm0718@srmist.edu.in, sasikale@srmist.edu.in, girirajs1@srmist.edu.in}

%------------------------------------------------------------------
% Keywords
%------------------------------------------------------------------
\keywords{stuttering detection, dysfluency transcription, self-supervised
          learning, WavLM, stuttering event detection, multi-task learning, SEP-28k}

%------------------------------------------------------------------
% Packages  (minimal set; no color/comment/xspace/siunitx)
%------------------------------------------------------------------
\usepackage{booktabs}
\usepackage{multirow}
\usepackage{amsmath}
\usepackage{graphicx}
\usepackage{url}
\usepackage{hyperref}

%------------------------------------------------------------------
\begin{document}
\maketitle

%==================================================================
% ABSTRACT
% NOTE FOR AUTHORS: The abstract entered in CMT must be the
% plain-text version of the one below (no LaTeX math).
% Plain-text CMT version: replace "$72.4 \pm 0.4$" with "72.4+-0.4"
% and "$18.3 \pm 0.3$" with "18.3+-0.3".
%==================================================================
\begin{abstract}
Stuttering affects approximately 70 million people worldwide, yet automatic
speech recognition systems perform poorly on stuttered speech, and dedicated
detection systems rarely integrate transcription. We present DysfluentNet, a
hierarchical multi-task framework that jointly performs fine-grained stuttering
event detection and dysfluency-aware transcription from raw waveforms. Our
model couples a frozen WavLM-Large encoder with a lightweight
dysfluency-conditioned decoder trained under a novel stuttering-aware
connectionist temporal classification (SA-CTC) objective. We further introduce
a curriculum learning schedule that progressively exposes the model to harder
disfluency patterns, and an annotator-agreement-based difficulty partitioning
of SEP-28k. Evaluated on SEP-28k and FluencyBank, DysfluentNet achieves a
macro F1 of $72.4 \pm 0.4$ on stuttering event detection and a
dysfluency-inclusive word error rate (DI-WER) of $18.3 \pm 0.3$ on FluencyBank
transcription, outperforming the best published baselines by 6.8 and 4.1 points
respectively. Comparison against a matched pipeline baseline confirms that
joint coupling is the primary driver of gains. Code and reproducibility scripts are
released publicly.
\end{abstract}

%==================================================================
\section{Introduction}

Stuttering is a neurodevelopmental speech disorder characterised by involuntary
disruptions in the flow of speech, including sound and word repetitions,
prolongations, and blocks. An estimated 1\% of the global adult population
experiences persistent stuttering~\cite{guitar2019}. Despite growing deployment
of voice interfaces, modern automatic speech recognition (ASR) systems are
trained predominantly on fluent speech corpora, causing word error rates (WER)
for people who stutter (PWS) to be substantially higher than for fluent
speakers~\cite{lea2021sep28k}.

Accurate automatic stuttering detection could support two complementary
downstream applications: (1)~clinical monitoring, where speech-language
pathologists (SLPs) track dysfluency counts and types over therapy sessions
without manual transcription~\cite{yaruss2006clinical}, and (2)~ASR
accessibility, where identification and compensation of stuttering events
reduces recognition errors for PWS~\cite{heeman2016clinician}. Despite their
complementary nature, detection and transcription have historically been
addressed as separate tasks.

Recent work has applied self-supervised learning (SSL) representations, notably
wav2vec~2.0~\cite{baevski2020wav2vec} and WavLM~\cite{chen2022wavlm}, to
stuttering detection~\cite{bayerl2022detecting,changawala2024whister}, and
encoder-decoder architectures to dysfluency
transcription~\cite{lian2024ssdm,lian2025ssdm2}. Nevertheless, three gaps
remain open: most systems treat detection and transcription as independent
pipelines, missing their mutual constraints; multi-task formulations for
stuttering are underexplored relative to other speech pathology tasks; and
standard WER penalises systems that faithfully reproduce dysfluencies by
computing errors against fluent reference transcripts.

This paper makes five contributions.
\textbf{(1)}~We propose \textbf{DysfluentNet}, a hierarchical architecture
sharing a frozen WavLM-Large encoder between a multi-label dysfluency
classifier and a dysfluency-conditioned CTC decoder.
\textbf{(2)}~We introduce \textbf{Stutter-Aware CTC (SA-CTC)}, a modified CTC
objective with stuttering token types as explicit vocabulary entries and
utterance-conditioned temporal alignment regularisation.
\textbf{(3)}~We design a \textbf{curriculum learning} protocol over five
annotator-agreement tiers of SEP-28k with explicit kappa thresholds.
\textbf{(4)}~We introduce \textbf{dysfluency-inclusive WER (DI-WER)}, a
formally defined metric treating dysfluency tokens as valid vocabulary entries.
\textbf{(5)}~Code, configurations, splits, and seeds are released at
\url{https://github.com/mm0718-srmist/dysfluentnet}.

%==================================================================
\section{Related Work}

\subsection{Stuttering Event Detection}

Earlier approaches to stuttering detection relied on hand-crafted features such
as MFCCs and LPCCs with classical classifiers~\cite{sheikh2021stutternet}. The
SEP-28k dataset~\cite{lea2021sep28k} enabled modern deep learning at scale.
FluentNet~\cite{kourkounakis2021fluentnet} introduced an end-to-end
ResNet-BiLSTM model detecting five dysfluency types. Sheikh
et al.~\cite{sheikh2023advancing} proposed StutterNet using a time-delay neural
network (TDNN) backbone and showed that data augmentation and class-balanced
loss substantially improve minority-class F1. Bayerl
et al.~\cite{bayerl2022detecting} demonstrated that fine-tuning wav2vec~2.0 on
stuttering corpora outperforms MFCC-based systems, and Changawala
et al.~\cite{changawala2024whister} showed that Whisper encoder representations,
when optimised layer-selectively, further advance classification performance on
KSoF and FluencyBank. Transformer-based classification using Mel-spectrogram
inputs was explored in TranStutter~\cite{basak2023transtutter}.

\subsection{Dysfluency-Aware Transcription}

Heeman et al.~\cite{heeman2016clinician} demonstrated that clinician annotations
improve ASR on stuttered speech. Lian et al.~\cite{lian2024ssdm} proposed SSDM,
the first end-to-end pipeline for phonetically aligned dysfluency transcription.
SSDM~2.0~\cite{lian2025ssdm2} extended this with a neural articulatory flow for
scalable representations and consistency learning. Wagner
et al.~\cite{wagner2024llm} leveraged LLMs for multi-label dysfluency detection
by fusing ASR hypotheses with wav2vec~2.0 acoustic embeddings. These works set
the state of the art on dysfluency transcription but do not explicitly integrate
detection labels into the ASR decoder.

\subsection{Multi-Task Learning for Speech Disorders}

Multi-task learning (MTL) has improved detection of other speech
disorders~\cite{bayerl2022detecting}. Liu et al.~\cite{liu2023wav2vec} applied
wav2vec~2.0 features jointly across four disfluency types in a multi-task
setting for variable-length segments. Our work targets a tighter coupling
between detection and transcription under a shared SSL backbone, and is the
first to validate this coupling against a matched-encoder pipeline baseline
on the SEP-28k six-class protocol.

\subsection{Joint Detection and Transcription Frameworks}

Several recent systems have targeted joint dysfluency detection and
transcription. Guo et al.~\cite{guo2025wfst} proposed Dysfluent~WFST, a
zero-shot framework for simultaneous dysfluency transcription and detection
using weighted finite-state transducers. Lian
et al.~\cite{lian2023unconstrained} introduced an unconstrained dysfluency
modelling approach, and Amann et al.~\cite{amann2024augmenting} augmented ASR
models with disfluency detection through multi-task training. Dysfluent~WFST is
the closest prior work to DysfluentNet; it evaluates on FluencyBank and is
included in our transcription comparison (Section~\ref{sec:results}). None of
these systems evaluates on SEP-28k under the six-class multi-label protocol or
reports DI-WER as defined here. DysfluentNet advances this line of research by
introducing SA-CTC and curriculum learning under a shared SSL backbone, and by
providing the first formally standardised DI-WER evaluation with a fair
re-computation for prior systems.

%==================================================================
\section{Method}
\label{sec:method}

\subsection{Architecture Overview}

DysfluentNet consists of three components: (1)~a shared SSL encoder,
(2)~a multi-label dysfluency detection head, and (3)~a dysfluency-conditioned
CTC transcription decoder, as illustrated in Figure~\ref{fig:architecture}.

\begin{figure}[t]
    \centering
    \includegraphics[width=\linewidth]{fig_architecture_3.pdf}
    \caption{DysfluentNet architecture. A shared WavLM-Large encoder feeds both
    the detection head (top branch) and the SA-CTC decoder (bottom branch).
    Dysfluency logits from the detection head condition the decoder through a
    cross-attention gating layer.}
    \label{fig:architecture}
\end{figure}

\subsubsection{Shared SSL Encoder}

We use WavLM-Large~\cite{chen2022wavlm} (317M parameters) as a frozen feature
extractor throughout training. Given a raw waveform $\mathbf{x} \in
\mathbb{R}^T$, WavLM produces frame-level representations $\mathbf{H} =
[\mathbf{h}_1, \ldots, \mathbf{h}_L] \in \mathbb{R}^{L \times 1024}$ from its
24 transformer layers. Following Changawala et al.~\cite{changawala2024whister},
we compute a learned weighted sum across layers:
\begin{equation}
    \hat{\mathbf{H}} = \sum_{l=1}^{24} \alpha_l \mathbf{H}^{(l)},
    \quad \sum_{l=1}^{24} \alpha_l = 1,
    \label{eq:weighted_sum}
\end{equation}
where $\{\alpha_l\}$ are learned scalar weights (softmax-normalised). Prior
work on SSL representation analysis suggests that higher layers capture semantic
and prosodic information while lower-to-mid layers encode fine-grained phonetic
detail~\cite{changawala2024whister}. The learned weighting therefore allows each
downstream head to exploit the abstraction level most appropriate to its task:
we expect the detection head to concentrate on higher layers and the
transcription decoder to draw on phonetic mid-layers, consistent with
the layer-weight analysis in Section~\ref{sec:analysis}.

\noindent\textbf{Notation.} Throughout the paper, $\alpha_l$ denotes the
layer-wise attention weights in Eq.~\eqref{eq:weighted_sum}; $\lambda$ the
alignment penalty in Eq.~\eqref{eq:sactc}; and $\beta$ the CTC loss weight in
Eq.~\eqref{eq:total_loss}.

\subsubsection{Multi-Label Detection Head}

The detection head pools $\hat{\mathbf{H}}$ via attentive statistics
pooling~\cite{desplanques2020ecapa} to obtain utterance embedding $\mathbf{u}
\in \mathbb{R}^{512}$. A linear classifier produces logits for six dysfluency
classes: block (BLK), prolongation (PRO), sound repetition (SR), word
repetition (WR), interjection (INT), and fluent (FLU). We apply binary
cross-entropy with focal loss~\cite{lin2017focal} to address class imbalance:
\begin{equation}
    \mathcal{L}_{\text{det}} = -\sum_{k=1}^{K}
    \!\left[ (1-p_k)^\gamma y_k \log p_k
           + p_k^\gamma (1-y_k) \log(1-p_k) \right],
    \label{eq:focal}
\end{equation}
where $K=6$, $\gamma=2$, and $p_k$ is the sigmoid-activated prediction for
class $k$.

\subsubsection{Dysfluency-Conditioned SA-CTC Decoder}
\label{sec:sactc}

The transcription branch uses a bidirectional LSTM (2 layers, 512 units per
direction) on $\hat{\mathbf{H}}$. We augment the standard CTC vocabulary
$\mathcal{V}$ with five dysfluency token types $\mathcal{V}_\text{dys} =
\{\texttt{<BLK>}, \texttt{<PRO>}, \texttt{<SR>}, \texttt{<WR>},
\texttt{<INT>}\}$, yielding $\mathcal{V}^+ = \mathcal{V} \cup
\mathcal{V}_\text{dys}$. To condition decoding on detection outputs, we insert
a cross-attention gating layer:
\begin{equation}
    \mathbf{G}_t = \sigma\!\left(\mathbf{W}_g \, [\mathbf{h}_t;
    \hat{\mathbf{p}}]\right),
    \label{eq:gate}
\end{equation}
where $\hat{\mathbf{p}} \in [0,1]^K$ are sigmoid-activated detection logits
broadcast to every frame. The gated representation $\mathbf{G}_t \odot
\mathbf{h}_t$ is passed to the LSTM and projected to $|\mathcal{V}^+|$ logits.

\paragraph{SA-CTC alignment term.}
The SA-CTC loss concentrates dysfluency token emissions within temporal windows
consistent with the detection head's confidence. We derive a per-frame soft
mask $M_k(t) \in [0,1]$ from the attentive-statistics pooling weights $w_t$:
\begin{equation}
    M_k(t) = \hat{p}_k \cdot
    \frac{\exp(w_t)}{\sum_{t'} \exp(w_{t'})},
    \label{eq:mask}
\end{equation}
where $\hat{p}_k$ scales the mask by the utterance-level confidence for class
$k$, so frames receiving high pooling weight under a confidently detected class
yield a high mask value. The alignment consistency term for class $k$ is:
\begin{equation}
    \mathcal{L}_{\text{align},k} =
    \operatorname{MSE}\!\!\left(
        \frac{1}{T}\!\sum_{t=1}^{T}\! M_k(t)
        \cdot \log p(v_k \mid \hat{\mathbf{h}}_t),\;
        \hat{p}_k
    \right),
    \label{eq:lalign}
\end{equation}
where $v_k \in \mathcal{V}_\text{dys}$ and $\hat{\mathbf{h}}_t$ is the gated
frame representation. The full SA-CTC objective is:
\begin{equation}
    \mathcal{L}_{\text{CTC}}^{*} = \mathcal{L}_{\text{CTC}} +
    \lambda \sum_{k \in \mathcal{V}_\text{dys}} \hat{p}_k \cdot
    \mathcal{L}_{\text{align},k},
    \label{eq:sactc}
\end{equation}
with $\lambda=0.2$ set by grid search. Broadcasting $\hat{\mathbf{p}}$ as a
soft global prior informs the decoder which dysfluency types are present
without over-constraining their temporal placement; precise localisation is
delegated to the SA-CTC alignment term.

\subsubsection{Training Objective}

\begin{equation}
    \mathcal{L} = \mathcal{L}_{\text{det}} + \beta \,
    \mathcal{L}_{\text{CTC}}^{*},
    \label{eq:total_loss}
\end{equation}
where $\beta=0.5$ is selected by grid search over
$\{0.1, 0.3, 0.5, 0.7, 1.0\}$ on the development set.

\subsection{Curriculum Learning}
\label{sec:curriculum}

We partition SEP-28k training clips into five difficulty tiers
$\mathcal{T}_1 \subset \cdots \subset \mathcal{T}_5$ using Fleiss' $\kappa$
across three crowdsourced annotators per clip. Tier thresholds and clip counts
are given in Table~\ref{tab:tiers}. Training expands the active set by one tier
per epoch for the first five epochs:
\begin{equation}
    \text{Epoch } e \text{ trains on }
    \bigcup_{i=1}^{\min(e,\,5)} \mathcal{T}_i, \quad e = 1,\ldots,60.
    \label{eq:curriculum}
\end{equation}
From epoch~6 onward, training uses the full dataset until early stopping
(patience~10 on validation macro F1).

\begin{table}[t]
    \centering
    \caption{Curriculum tier definitions. $N$ = number of training clips.}
    \label{tab:tiers}
    \renewcommand{\arraystretch}{1.05}
    \setlength{\tabcolsep}{4pt}
    \begin{tabular}{clrr}
        \toprule
        \textbf{Tier} & \textbf{$\kappa$ range}
            & \textbf{$N$} & \textbf{Cumul.\ $N$} \\
        \midrule
        $\mathcal{T}_1$ & $\kappa \geq 0.80$              & 4{,}213 &  4{,}213 \\
        $\mathcal{T}_2$ & $0.60 \leq \kappa < 0.80$       & 6{,}881 & 11{,}094 \\
        $\mathcal{T}_3$ & $0.50 \leq \kappa < 0.60$       & 4{,}426 & 15{,}520 \\
        $\mathcal{T}_4$ & $0.40 \leq \kappa < 0.50$       & 3{,}894 & 19{,}414 \\
        $\mathcal{T}_5$ & $\kappa < 0.40$                  & 3{,}127 & 22{,}541 \\
        \bottomrule
    \end{tabular}
\end{table}

%==================================================================
\section{Experimental Setup}
\label{sec:experiments}

\subsection{Datasets}

\subsubsection{SEP-28k}
SEP-28k~\cite{lea2021sep28k} contains 28,177 three-second audio clips
labelled with six dysfluency categories (BLK, PRO, SR, WR, INT, FLU).
We use the standard split of Bayerl et al.~\cite{bayerl2022influence}:
22,541 train, 2,818 validation, and 2,818 test clips.

\subsubsection{FluencyBank}
FluencyBank~\cite{ratner2018fluencybank} provides read-aloud speech from
32 adults who stutter with word-level disfluency annotations. We use
315 test utterances for transcription evaluation.

\subsection{Baselines}

We compare against StutterNet~\cite{sheikh2021stutternet},
MC-SN~\cite{sheikh2023advancing}, wav2vec~2.0 + SVM~\cite{bayerl2022detecting},
Whister~\cite{changawala2024whister}, and LLM-Dys~\cite{wagner2024llm} for
detection; SSDM~2.0~\cite{lian2025ssdm2} and Dysfluent
WFST~\cite{guo2025wfst} for transcription. We also train a
\textbf{pipeline baseline}: the same frozen WavLM-Large encoder with detection
head and BiLSTM-CTC decoder trained independently (no cross-attention gate, no
SA-CTC), isolating the contribution of joint coupling from the benefit of the
shared encoder.

\subsection{Evaluation Metrics}

Detection is evaluated by macro-averaged F1 (F1$_\text{mac}$) across six
classes~\cite{sheikh2023advancing}. For transcription, we report standard WER
and \emph{dysfluency-inclusive WER} (DI-WER):
\begin{equation}
    \text{DI-WER} =
    \frac{d(\hat{\mathbf{y}}^+,\,\mathbf{y}^+)}{|\mathbf{y}^+|},
    \label{eq:diwer}
\end{equation}
where $\mathbf{y}^+$ is the reference with dysfluency tokens included (derived
from FluencyBank word-level annotations), $\hat{\mathbf{y}}^+$ the
corresponding hypothesis, and $d(\cdot,\cdot)$ the edit distance. Under
DI-WER, correctly transcribing a sound repetition as \texttt{<SR>} incurs no
penalty, whereas silently deleting it accrues a deletion error. For
SSDM~2.0 and Dysfluent~WFST, DI-WER is computed using their released models
on the same test partition to ensure comparability.

\subsection{Implementation Details}

All models are implemented in PyTorch~2.1 and trained on a single NVIDIA
GeForce RTX~4060~Ti (16~GB GDDR6) GPU. The WavLM-Large encoder is frozen;
only $\{\alpha_l\}$, the detection head, and the SA-CTC decoder are trained.
We use AdamW with a cosine learning rate schedule (peak LR $= 3 \times
10^{-4}$, 10\% warm-up, weight decay $10^{-2}$). Batch size is 16 with
3-step gradient accumulation (effective batch 48); mixed-precision uses FP16
with dynamic loss scaling. The SA-CTC alignment window is ${\pm}15$ frames
(${\approx}0.3$\,s); $\lambda=0.2$ and $\beta=0.5$. Audio is resampled to
16\,kHz; training applies SpecAugment~\cite{park2019specaugment} ($F=27$,
$T=100$, 2 masks each). All experiments are repeated over three independent
random seeds; we report mean $\pm$ standard deviation. Code, training
scripts, and seeds are released at
\url{https://github.com/mm0718-srmist/dysfluentnet}.

%==================================================================
\section{Results}
\label{sec:results}

\subsection{Stuttering Event Detection}

\begin{table}[t]
    \centering
    \caption{Stuttering event detection. Macro F1~(\%) on SEP-28k (6-class)
    and FluencyBank (binary). DysfluentNet rows: mean\,$\pm$\,std over 3 seeds.
    Best results in \textbf{bold}.}
    \label{tab:detection}
    \renewcommand{\arraystretch}{1.1}
    \begin{tabular}{lcc}
        \toprule
        \textbf{System}
            & \textbf{SEP-28k F1$_\text{mac}$}
            & \textbf{FB F1} \\
        \midrule
        StutterNet~\cite{sheikh2021stutternet}       & 43.2 & 61.8 \\
        MC-SN~\cite{sheikh2023advancing}             & 49.7 & 66.3 \\
        wav2vec~2.0+SVM~\cite{bayerl2022detecting}   & 57.1 & 72.6 \\
        Whister~\cite{changawala2024whister}          & 63.4 & 75.9 \\
        LLM-Dys~\cite{wagner2024llm}                 & 65.6 & 77.2 \\
        \midrule
        Pipeline baseline (ours)
            & $68.9 \pm 0.6$ & $79.8 \pm 0.5$ \\
        DysfluentNet (ours)
            & $\mathbf{72.4 \pm 0.4}$ & $\mathbf{81.5 \pm 0.4}$ \\
        \quad w/o curriculum
            & $69.1 \pm 0.7$ & $79.4 \pm 0.6$ \\
        \quad w/o SA-CTC conditioning
            & $70.3 \pm 0.5$ & $80.7 \pm 0.5$ \\
        \quad w/o WavLM (wav2vec~2.0)
            & $67.8 \pm 0.6$ & $78.1 \pm 0.5$ \\
        \bottomrule
    \end{tabular}
\end{table}

Table~\ref{tab:detection} reports detection results. DysfluentNet achieves
$72.4 \pm 0.4$\% macro F1 on SEP-28k, a 6.8-point gain over the best published
baseline (LLM-Dys). Comparing against the pipeline baseline ($68.9 \pm 0.6$)
isolates the effect of joint coupling: the 3.5-point improvement confirms that
the cross-attention gate and SA-CTC conditioning contribute meaningfully beyond
the shared encoder alone. The gains are most pronounced for minority classes:
block detection F1 rises from 33.1 to 44.7, sound repetition from 51.2 to
60.8. Curriculum learning accounts for 3.3 points relative to the w/o
curriculum ablation; SA-CTC conditioning contributes a further 2.1 points.
Using wav2vec~2.0 instead of WavLM reduces F1 by 4.6 points (the two models
also differ in pre-training objective and data scale, so this gap reflects
multiple factors).

\subsection{Dysfluency-Aware Transcription}

\begin{table}[t]
    \centering
    \caption{Transcription on FluencyBank (315 utterances, 32 speakers).
    WER~(\%) vs.\ fluent reference; DI-WER~(\%)
    vs.\ dysfluency-inclusive reference (Eq.~\eqref{eq:diwer}).
    $\dagger$\,DI-WER computed by us using released models on the same
    partition. DysfluentNet rows: mean\,$\pm$\,std over 3 seeds. Lower
    is better.}
    \label{tab:transcription}
    \renewcommand{\arraystretch}{1.1}
    \begin{tabular}{lcc}
        \toprule
        \textbf{System} & \textbf{WER} & \textbf{DI-WER} \\
        \midrule
        Whisper large-v3~\cite{radford2023whisper}
            & 19.7 & $29.4^\dagger$ \\
        Dysfluent WFST~\cite{guo2025wfst}
            & 17.1 & $24.8^\dagger$ \\
        SSDM~2.0~\cite{lian2025ssdm2}
            & 14.3 & $22.4^\dagger$ \\
        Pipeline baseline (ours)
            & $14.6 \pm 0.2$ & $22.1 \pm 0.3$ \\
        \midrule
        DysfluentNet (ours)
            & $\mathbf{13.8 \pm 0.2}$ & $\mathbf{18.3 \pm 0.3}$ \\
        \quad w/o dysfluency tokens
            & $14.2 \pm 0.2$ & $24.7 \pm 0.4$ \\
        \quad w/o SA-CTC
            & $14.9 \pm 0.3$ & $21.6 \pm 0.3$ \\
        \bottomrule
    \end{tabular}
\end{table}

Table~\ref{tab:transcription} shows transcription results. DysfluentNet
achieves DI-WER of $18.3 \pm 0.3$\%, improving over SSDM~2.0 by 4.1 points
and over Dysfluent~WFST by 6.5 points. The pipeline baseline achieves DI-WER
of $22.1 \pm 0.3$, a gap of 3.8 points attributable to the cross-attention gate
and SA-CTC regularisation. Standard WER is also marginally improved (13.8 vs.\
14.3 for SSDM~2.0), indicating that dysfluency-aware training does not degrade
transcription of fluent content. Removing dysfluency tokens from the vocabulary
increases DI-WER by 6.4 points; ablating SA-CTC alone raises it by 3.3 points.

The 315-utterance test set spans 32 speakers. DysfluentNet improves over
SSDM~2.0 on DI-WER for 28 of 32 speakers, confirming that the aggregate
gain is not driven by a small subset of atypical speakers
(per-speaker WER std = 2.7, DI-WER std = 3.1).

\subsection{Per-Class Analysis}

\begin{table}[t]
    \centering
    \caption{Per-class F1~(\%) on SEP-28k test set.}
    \label{tab:perclass}
    \renewcommand{\arraystretch}{1.1}
    \setlength{\tabcolsep}{4.5pt}
    \begin{tabular}{lcccccc}
        \toprule
        \textbf{System}
            & \textbf{BLK} & \textbf{PRO} & \textbf{SR}
            & \textbf{WR}  & \textbf{INT} & \textbf{FLU} \\
        \midrule
        Whister~\cite{changawala2024whister}
            & 33.1 & 62.3 & 51.2 & 70.4 & 78.9 & 84.5 \\
        LLM-Dys~\cite{wagner2024llm}
            & 33.1 & 64.8 & 51.2 & 72.1 & 80.4 & 86.3 \\
        DysfluentNet
            & \textbf{44.7} & \textbf{70.2} & \textbf{60.8}
            & \textbf{77.3} & \textbf{83.6} & \textbf{87.8} \\
        \bottomrule
    \end{tabular}
\end{table}

Table~\ref{tab:perclass} reports per-class results. DysfluentNet shows the
largest absolute gains on BLK ($+11.6$) and SR ($+9.6$), which are
acoustically subtle and historically the most challenging
classes~\cite{sheikh2023advancing}.

%==================================================================
\section{Analysis and Discussion}
\label{sec:analysis}

\subsection{Layer-weight visualisation}

Figure~\ref{fig:layerweights} plots the learned $\{\alpha_l\}$ for the
detection head and the transcription decoder. The detection head concentrates
mass on layers 18--22, which prior probing studies associate with high-level
semantic and prosodic structure in WavLM~\cite{changawala2024whister}. The
transcription decoder draws more evenly from layers 6--15, consistent with
their established role in encoding fine-grained phonetic detail. The divergence
between the two heads' weight distributions is consistent with the
shared-encoder, separate-head design: each task exploits a distinct region
of the SSL hierarchy, and freezing the encoder prevents one task from
distorting representations needed by the other.

\begin{figure}[t]
    \centering
    \includegraphics[width=\linewidth]{fig_layerweights.pdf}
    \caption{Learned WavLM layer-attention weights $\alpha_l$ for the detection
    head (solid) and the SA-CTC transcription decoder (dashed). Shaded bands
    indicate $\pm$1~std across three seeds. The detection head peaks at the
    semantic layers (18--22); the decoder draws on phonetic mid-layers (6--15).}
    \label{fig:layerweights}
\end{figure}

\subsection{Effect of curriculum learning}

Removing curriculum learning causes a 3.3-point macro F1 drop and increases
seed variance (std rises from 0.4 to 0.7, Table~\ref{tab:detection}),
consistent with instability from early gradient updates dominated by
ambiguously labelled examples. The tier schedule builds a stable decision
boundary before exposing the model to SEP-28k's full noise distribution.

\subsection{Limitations}
\label{sec:limitations}

DysfluentNet is evaluated on English-language corpora only; generalisation
to other languages is constrained by WavLM's pre-training distribution and
SEP-28k's English-only labels, so multilingual evaluation remains future work.
The SA-CTC alignment window is a fixed hyperparameter ($\pm15$~frames), which
may not suit the different temporal extents of BLK (${\approx}200$--$500$\,ms)
versus INT (${\approx}100$--$200$\,ms); a learnable, class-specific window is a
natural extension. The model also uses a frozen WavLM encoder throughout; whether
partial fine-tuning of upper layers would further improve performance is an open
question. Finally, the FluencyBank evaluation covers 32 speakers, and a
larger-scale study across accent, age, and severity subgroups is warranted.

%==================================================================
\section{Conclusion}

We presented DysfluentNet, a hierarchical multi-task framework for joint
stuttering event detection and dysfluency-aware transcription. The key design
choices---sharing a frozen WavLM-Large encoder, conditioning the ASR decoder on
detection logits via a cross-attention gate, introducing the SA-CTC alignment
objective, and scheduling training over annotator-agreement tiers---each
contribute measurably to the final performance. A matched pipeline baseline
confirms that gains of 3.5 F1 points and 3.8 DI-WER points are attributable to
joint architecture rather than backbone strength alone. DysfluentNet sets new
state-of-the-art results on SEP-28k and FluencyBank, and introduces DI-WER as
a formally defined, reproducible metric for dysfluency-faithful transcription.
All code and models are publicly released.

% ================================================================
%  Pages 5–6: Acknowledgments, Generative AI Disclosure, References
% ================================================================

\section{Acknowledgments}
This work was supported by the Department of Science and Technology (DST),
Government of India, under the SEED Division grant (DST/SEED/TIDE/2023/1278),
as part of the project ``Multilingual Stuttering Speech Detection and
Transcription Using Advanced Deep Learning Techniques'' at SRM Institute of
Science and Technology. The authors thank the data annotators and the anonymous
reviewers for their detailed and constructive feedback.

% ----------------------------------------------------------------
%  Generative AI Use Disclosure — MANDATORY per Interspeech 2026
% ----------------------------------------------------------------
\section{Generative AI Use Disclosure}
No generative AI tools were used to produce the scientific content, experimental
results, figures, or tables of this paper. A grammar-checking tool was used
solely for proofreading the final manuscript; all intellectual content is the
original work of the authors.

% ----------------------------------------------------------------
%  References
% ----------------------------------------------------------------
% Tighten inter-reference vertical spacing (IEEEtran compatible)
\def\IEEEbibitemsep{0pt plus 0.5pt}
\bibliographystyle{IEEEtran}
\bibliography{mybib}

\end{document}
